\section{Differentiable Simulation}
\label{sec:diffsim}

A physics simulator is usually treated as a black box: we give it an
initial state and a sequence of control inputs, and it tells us what
happens. A \emph{differentiable} simulator also tells us how the outcome
would change if we changed the inputs a little. Formally, it returns the
derivative of the outcome with respect to the inputs. This chapter
introduces the idea: what such a derivative is, why it is useful, how it is
computed, and why it is harder to get right than it first looks, especially
when objects touch each other.

\subsection{What Is a Differentiable Simulator?}
\label{sec:diffsim-what}

\paragraph{A simulator as a function.}
A simulator advances a state $\mathbf{s}_k$ (positions and velocities of
all bodies) by one time step of length $h$:
\begin{equation}
  \mathbf{s}_{k+1} = \mathbf{f}(\mathbf{s}_k, \mathbf{a}_k, \boldsymbol{\theta}),
  \qquad k = 0, \dots, T-1 .
  \label{eq:diffsim-step}
\end{equation}
Here $\mathbf{a}_k$ are the control inputs at step $k$, for example motor
torques or target wheel velocities, and $\boldsymbol{\theta}$ are the
parameters of the model, for example masses, friction coefficients or the
lengths of the links of a robot. Running $T$ steps from an initial state
$\mathbf{s}_0$ gives a trajectory $\mathbf{s}_1, \dots, \mathbf{s}_T$.

What we care about is usually summarized in one number, the \emph{loss}
\begin{equation}
  L = \sum_{k=1}^{T} \ell_k(\mathbf{s}_k, \mathbf{a}_{k-1}) ,
  \label{eq:diffsim-loss}
\end{equation}
for example the distance between the final position of a robot and a
goal, plus a penalty on the control effort. Because every state depends on
the inputs, $L$ is a function of the initial state, the controls and the
parameters. A differentiable simulator computes, in addition to $L$, its
gradient
\begin{equation}
  \nabla L = \left(
    \frac{\partial L}{\partial \mathbf{s}_0},\;
    \frac{\partial L}{\partial \mathbf{a}_0}, \dots,
    \frac{\partial L}{\partial \mathbf{a}_{T-1}},\;
    \frac{\partial L}{\partial \boldsymbol{\theta}}
  \right) .
\end{equation}
The gradient points in the direction in which the loss increases fastest.
Stepping against it lowers the loss, which is the basic step of every
gradient-based optimizer.

\paragraph{What it is used for.}
With the gradient available, many problems in robotics and graphics become
optimization problems that standard gradient methods can solve
\citep{newbury2024review}:
\begin{itemize}
  \item \textbf{Trajectory optimization.} Find the controls
    $\mathbf{a}_0, \dots, \mathbf{a}_{T-1}$ that make a robot reach a
    goal, throw a ball into a basket or walk without falling. The unknowns
    are the controls.
  \item \textbf{System identification.} Find the parameters
    $\boldsymbol{\theta}$ (a mass, a friction coefficient) for which the
    simulated trajectory matches a recorded real one. The loss is the
    distance between the two trajectories.
  \item \textbf{Policy learning.} Train a neural network that maps the
    state to the controls, $\mathbf{a}_k = \boldsymbol{\pi}_{\mathbf{w}}(\mathbf{s}_k)$.
    The gradient of the loss with respect to the network weights
    $\mathbf{w}$ flows backwards through the simulator into the network,
    in the same way as through any other layer.
  \item \textbf{Design optimization.} Change the shape or the dimensions
    of a mechanism so that it performs a task better.
\end{itemize}

\subsection{Computing the Gradient}
\label{sec:diffsim-gradient}

There are three basic ways to compute $\nabla L$. Table~\ref{tab:diffsim-methods}
compares them at the end of this section.

\paragraph{Finite differences.}
The simplest way is to perturb one input $p_i$ by a small $\varepsilon$
and simulate again:
\begin{equation}
  \frac{\partial L}{\partial p_i}
  \approx \frac{L(\mathbf{p} + \varepsilon \mathbf{e}_i) - L(\mathbf{p} - \varepsilon \mathbf{e}_i)}{2\varepsilon} .
  \label{eq:diffsim-fd}
\end{equation}
This \emph{central difference} needs two simulations per input, so the full
gradient of $P$ inputs costs $2P$ simulations. Its accuracy is limited from
two sides. A large $\varepsilon$ gives a truncation error of order
$\varepsilon^2$, because the formula assumes $L$ is a parabola near
$\mathbf{p}$. A small $\varepsilon$ makes the two values of $L$ nearly
equal, and their difference loses the significant digits of the rounding
error $\delta$ of $L$, an error of order $\delta/\varepsilon$. The best
$\varepsilon$ balances the two, and in single precision the result rarely
has more than three or four correct digits
\citep{griewank2008evaluating}. Finite differences are therefore too
expensive and too inaccurate for optimization with many unknowns, but they
are the standard way to \emph{check} a gradient computed by other means
(Section~\ref{sec:diffsim-quality}).

\paragraph{Automatic differentiation.}
A simulator is a long sequence of elementary operations (additions,
multiplications, square roots, \dots), each of which has a known
derivative. \emph{Automatic differentiation} (AD) applies the chain rule to
this sequence mechanically and gives derivatives that are exact up to
rounding \citep{griewank2008evaluating,baydin2018automatic}. It comes in
two modes.
\begin{itemize}
  \item \textbf{Forward mode} carries, next to every intermediate value,
    its derivative with respect to \emph{one} chosen input. One pass gives
    the derivative of all outputs with respect to that input, so the full
    gradient of $P$ inputs needs $P$ passes, as with finite differences,
    but without their error.
  \item \textbf{Reverse mode} first runs the computation forward and
    records it (the \emph{tape}). It then goes through the record
    backwards and propagates the derivative of \emph{one} output, the loss,
    to all intermediate values and finally to all inputs. One backward
    pass gives the whole gradient, at a cost of a small multiple of the
    forward computation, independently of $P$. Backpropagation in neural
    networks is reverse-mode AD.
\end{itemize}
Since a loss is one number and the inputs are many, reverse mode is the
natural choice for simulation. Its price is memory: the backward pass needs
the intermediate values of the forward pass, so the whole trajectory must
be stored. For long trajectories, \emph{checkpointing} stores only some of
the states and recomputes the rest during the backward pass, trading
computation for memory \citep{griewank2008evaluating}. Several simulation
frameworks provide reverse-mode AD for their kernels, for example
DiffTaichi \citep{hu2020difftaichi}, Brax \citep{freeman2021brax} and
NVIDIA Warp.

\paragraph{Backpropagation through time.}
Applied to the trajectory \eqref{eq:diffsim-step}, reverse mode becomes a
backward recursion over the time steps. Let
$\bar{\mathbf{s}}_k = \partial L / \partial \mathbf{s}_k$ denote the
\emph{adjoint} of the state, i.e.\ the sensitivity of the loss to the state
at step $k$, including the influence of $\mathbf{s}_k$ on all later
states. Starting from $\bar{\mathbf{s}}_T = \partial \ell_T / \partial \mathbf{s}_T$,
the chain rule gives
\begin{equation}
  \bar{\mathbf{s}}_k
  = \frac{\partial \ell_k}{\partial \mathbf{s}_k}
  + \left(\frac{\partial \mathbf{f}}{\partial \mathbf{s}_k}\right)^{\!\top} \bar{\mathbf{s}}_{k+1},
  \qquad
  \frac{\partial L}{\partial \mathbf{a}_k}
  = \frac{\partial \ell_{k+1}}{\partial \mathbf{a}_k}
  + \left(\frac{\partial \mathbf{f}}{\partial \mathbf{a}_k}\right)^{\!\top} \bar{\mathbf{s}}_{k+1},
  \label{eq:diffsim-bptt}
\end{equation}
for $k = T-1, \dots, 0$, and the parameter gradient
$\partial L/\partial \boldsymbol{\theta}$ is the sum of
$(\partial \mathbf{f}/\partial \boldsymbol{\theta})^\top \bar{\mathbf{s}}_{k+1}$
over all steps. Each backward step multiplies the adjoint by the
transposed \emph{step Jacobian} $\partial \mathbf{f}/\partial \mathbf{s}_k$.
The matrix itself is never formed: what is needed is only its product with
a vector, which is what reverse-mode AD of one time step provides.

\paragraph{Implicit time stepping and the implicit function theorem.}
Many simulators, especially those for stiff systems and contact, do not
compute $\mathbf{s}_{k+1}$ by an explicit formula. They define it
\emph{implicitly}, as the solution of a system of nonlinear equations
\begin{equation}
  \mathbf{R}(\mathbf{s}_{k+1}, \mathbf{s}_k, \mathbf{a}_k, \boldsymbol{\theta}) = \mathbf{0},
  \label{eq:diffsim-residual}
\end{equation}
which they solve by an iterative method, typically Newton's method.
Implicit integrators are popular because they stay stable with large time
steps. Differentiating through the Newton iterations with AD is possible
but wasteful, since the iterations are only the road to the solution, and
the derivative of the road is not what we want.

The \emph{implicit function theorem} gives the derivative of the solution
directly. If $\mathbf{R} = \mathbf{0}$ holds at the solution and we change
$\mathbf{s}_k$ by a small $\mathrm{d}\mathbf{s}_k$, then $\mathbf{s}_{k+1}$
must change by a $\mathrm{d}\mathbf{s}_{k+1}$ that keeps $\mathbf{R}$ at
zero to first order:
\begin{equation*}
  \frac{\partial \mathbf{R}}{\partial \mathbf{s}_{k+1}}\,\mathrm{d}\mathbf{s}_{k+1}
  + \frac{\partial \mathbf{R}}{\partial \mathbf{s}_k}\,\mathrm{d}\mathbf{s}_k = \mathbf{0}
  \quad\Longrightarrow\quad
  \frac{\partial \mathbf{s}_{k+1}}{\partial \mathbf{s}_k}
  = -\left(\frac{\partial \mathbf{R}}{\partial \mathbf{s}_{k+1}}\right)^{\!-1}
     \frac{\partial \mathbf{R}}{\partial \mathbf{s}_k} ,
\end{equation*}
provided that $\partial \mathbf{R}/\partial \mathbf{s}_{k+1}$, the matrix
of the last Newton iteration, is invertible. Inserting the transpose of
this Jacobian into \eqref{eq:diffsim-bptt} gives the \emph{adjoint method}:
each backward step first solves one linear system with the transposed
Newton matrix and then multiplies by $\partial \mathbf{R}/\partial \mathbf{s}_k$,
\begin{equation}
  \left(\frac{\partial \mathbf{R}}{\partial \mathbf{s}_{k+1}}\right)^{\!\top} \boldsymbol{\mu}_{k+1}
  = \bar{\mathbf{s}}_{k+1},
  \qquad
  \bar{\mathbf{s}}_k
  = \frac{\partial \ell_k}{\partial \mathbf{s}_k}
  - \left(\frac{\partial \mathbf{R}}{\partial \mathbf{s}_k}\right)^{\!\top} \boldsymbol{\mu}_{k+1} .
  \label{eq:diffsim-adjoint}
\end{equation}
The gradients with respect to $\mathbf{a}_k$ and $\boldsymbol{\theta}$
follow in the same way, with $\partial \mathbf{R}/\partial \mathbf{a}_k$
and $\partial \mathbf{R}/\partial \boldsymbol{\theta}$ in place of
$\partial \mathbf{R}/\partial \mathbf{s}_k$, and they reuse the same
$\boldsymbol{\mu}_{k+1}$. The backward step therefore costs one linear
solve, no matter how many Newton iterations the forward step needed, and
it needs to store only the converged states. The same idea underlies the
adjoint-state method in engineering and geophysics \citep{plessix2006review},
neural ordinary differential equations \citep{chen2018neural} and
optimization layers in neural networks \citep{amos2017optnet}. Its
assumptions also mark its limits. The theorem describes the
\emph{exact} solution of \eqref{eq:diffsim-residual}, so if Newton's
method stopped early, or the linear solve in \eqref{eq:diffsim-adjoint}
is inexact, the computed gradient is only an approximation.

\begin{table}[h]
  \centering
  \small
  \begin{tabular}{@{}p{3.2cm}p{3.4cm}p{3.0cm}p{4.6cm}@{}}
    \toprule
    Method & Cost of $\nabla L$ & Memory & Accuracy \\
    \midrule
    Finite differences & $2P$ simulations & one trajectory & a few digits; step-size dependent \\
    Forward-mode AD & $P$ passes & one trajectory & exact up to rounding \\
    Reverse-mode AD & $\approx 1$ backward pass & whole trajectory and tape & exact up to rounding \\
    Implicit adjoint & one linear solve per step & converged states & exact if forward and adjoint solves are exact \\
    \bottomrule
  \end{tabular}
  \caption{Ways to compute the gradient of a scalar loss with respect to
    $P$ inputs. ``Exact'' refers to the derivative of the computation that
    was actually performed, which is not the same as the derivative of the
    physics (Section~\ref{sec:diffsim-contact}).}
  \label{tab:diffsim-methods}
\end{table}

\subsection{Why Contact Is Hard}
\label{sec:diffsim-contact}

So far we assumed that $\mathbf{f}$ is a smooth function. Free flight,
springs and joints satisfy this. Contact does not, and contact is what
robots use to walk, drive and grasp.

\paragraph{Contact makes the dynamics non-smooth.}
Drop a ball from height $y_0$ onto the floor and record its height after a
fixed time. While the ball is in the air, the final height depends
smoothly on $y_0$. Raise $y_0$ slightly, and the moment of impact moves
past the end of the simulation, so the ball suddenly ends up much higher
than the neighbouring trajectories that did bounce. The final height has a
\emph{kink} (its derivative jumps) or even a \emph{jump} as a function of
$y_0$. Contact introduces such features in three ways:
\begin{itemize}
  \item \textbf{Making or breaking contact.} A non-penetration constraint
    is active only when the bodies touch. Whether it is active changes
    abruptly.
  \item \textbf{Impacts.} A collision changes the velocity in an instant.
    Small changes of the initial state can change the order or the number
    of impacts.
  \item \textbf{Stick and slip.} Coulomb friction holds a body in place
    until the tangential force reaches $\mu$ times the normal force, and
    then lets it slide. The switch between sticking and sliding is again
    abrupt.
\end{itemize}

\paragraph{What the gradient sees.}
A derivative describes the function in an infinitesimal neighbourhood. At
a kink or a jump it either does not exist or, where it exists, it tells
nothing about what lies a small step away. Two failures are typical.
First, \emph{flat regions}: a box resting on a table under static friction
does not move when we push it slightly, so the derivative of its final
position with respect to the push is exactly zero, and gradient descent
never learns that a harder push would move it. Second, \emph{misleading
directions}: near a jump the gradient describes only one side of it, and a
step in its direction can increase the loss. A gradient that is correct
in the mathematical sense can still be useless for optimization, and
\citet{suh2022differentiable} show examples where it makes the optimizer
worse than a method that only samples the loss at random inputs. \citet{zhong2022differentiable} compare
four common contact models and find that their gradients are not always
correct, even in simple scenes.

\paragraph{Approaches.}
Differentiable simulators handle contact in roughly three ways.
\begin{itemize}
  \item \textbf{Analytic gradients of a hard contact model.} The contact
    forces are the solution of a complementarity problem or of a convex
    optimization problem, and the gradient follows from differentiating
    its optimality conditions, an application of the implicit function
    theorem above \citep{belbuteperes2018endtoend,werling2021fast,howell2022dojo}.
    The gradients are exact for the model, including the kinks.
  \item \textbf{Smoothing the model.} Compliant contact replaces the hard
    constraint by a stiff spring, and soft or relaxed versions of the
    complementarity conditions round off the kinks \citep{geilinger2020add}.
    The gradients become informative around the contact event, at the
    price of a less accurate physical model, and the stiffness controls
    the trade-off.
  \item \textbf{Smoothing the gradient.} Instead of changing the model,
    one averages gradients or function values over randomly perturbed
    inputs. The average is the gradient of a smoothed loss, which has no
    kinks. It is more expensive and noisier, but it can step over the
    flat regions \citep{suh2022differentiable}.
\end{itemize}
In all three cases the gradient is the gradient of \emph{some} smooth
problem close to the real one. Which one is best depends on the task, and
it is an active research topic.

\subsection{Gradient Quality}
\label{sec:diffsim-quality}

Even when the model is smooth, the computed gradient can be wrong in ways
that are not obvious from the code. It is useful to separate three
sources of error.

\paragraph{The wrong derivative.}
The backward pass must differentiate exactly the computation the forward
pass performed. Any mismatch, for example a forward step that uses one
value of a parameter while the backward step linearizes around another,
gives a gradient that is self-consistent but belongs to a different
problem. Such bugs do not cause crashes or warnings; they only make the
optimization worse.

\paragraph{Inexact solves.}
Implicit simulators stop Newton's method and the linear solver at a
tolerance, and so the implicit function theorem is applied at a point
that is not exactly a solution, and the linear system of
\eqref{eq:diffsim-adjoint} is solved only approximately. How much error
this introduces depends on the conditioning of the matrices involved, and
also on how the solution of the linear system is turned into the
gradient: a formula that subtracts two nearly equal quantities amplifies
a small relative error of the solution into a large error of the
gradient. Studies in computational fluid dynamics have found that
truncated adjoint solves can produce gradients that point in a wrong
direction \citep{akbarzadeh2020consistency}.

\paragraph{Long horizons and chaos.}
By \eqref{eq:diffsim-bptt}, the adjoint of the initial state is a product
of $T$ transposed step Jacobians applied to the final adjoint. If these
Jacobians tend to stretch vectors, the gradient grows exponentially with
$T$ (\emph{exploding gradients}); if they tend to shrink them, the
gradient vanishes. This is the same problem as in recurrent neural
networks. In chaotic systems, where nearby trajectories diverge
exponentially, the exact gradient of a long rollout is huge and its
direction carries little information about the loss at a finite distance
\citep{metz2021gradients}. Practical remedies are short horizons,
splitting the trajectory into shorter segments that are optimized
together (multiple shooting), and normalizing or clipping the gradient.

\paragraph{Checking a gradient.}
The standard test compares the computed gradient with finite differences
\eqref{eq:diffsim-fd}. Instead of all $P$ components, one usually checks a
few random directions $\mathbf{d}$: the directional derivative
$\nabla L^\top \mathbf{d}$ should agree with
$\left(L(\mathbf{p} + \varepsilon\mathbf{d}) - L(\mathbf{p} - \varepsilon\mathbf{d})\right)/(2\varepsilon)$.
Agreement is reported as a relative error, and for whole gradients as the
\emph{cosine} between the computed and the reference gradient: a cosine
near $1$ means the direction is right, which is what a gradient-based
optimizer needs most. Because finite differences are themselves
inaccurate, such checks are best done in double precision and with a few
values of $\varepsilon$.

\paragraph{Summary.}
A differentiable simulator turns a simulation into a function whose
gradient can be computed at roughly the cost of one more simulation, by
reverse-mode differentiation and, for implicit time stepping, by one
linear solve per step with the adjoint method. This makes optimization
over many unknowns feasible. The gradient is, however, only as good as the
model it differentiates and the numerics that compute it: contact makes
the true dynamics non-smooth, inexact solves and unstable formulas
introduce numerical errors, and long horizons amplify both. Using a
differentiable simulator well therefore means not only computing the
gradient, but also knowing how far it can be trusted.
